\title{Structural Relaxation in Glass-Forming Liquids at Extreme Pressure: Photon-Correlation Spectroscopy and Constrained Thermodynamic Surface Modeling}

\author{Kevin Wayne Lyon}
%\degreeyear/{May 2017}
%\degreeyear{July 2017}

\degreetitle{Doctor of Philosophy in Physics} 

\field{Field of study}
\chair{Dr. Gea Banacloche}

\othermembers{
Professor 1, Committee member\\ 
Professor 2, Committee member\\
Professor 3, Committee member\\
Professor 4, Committee member\\
}
\numberofmembers{5} % |chair| + |cochair| + |othermembers|
%%%%%END stuff that does not work
%%%%%%%

\abstract{
		Structural relaxation in glass-forming liquids is strongly dependent on both temperature and pressure, yet direct measurement of structural $\alpha$-relaxation under multi-gigapascal compression has remained experimentally challenging. This work extends depolarized photon correlation spectroscopy (DPCS) to pressures approaching $5\,\mathrm{GPa}$ in a diamond anvil cell, representing the highest-pressure implementation of this technique and enabling direct measurement of the structural $\alpha$-relaxation in a pure liquid at extreme compression. This capability provides access to dynamical regimes previously unexplored by photon correlation spectroscopy and establishes a platform for probing glassy dynamics at high pressure. To describe these measurements, a constrained relaxation-time surface framework is developed within the Vogel--Tammann--Fulcher (VTF) formalism. A polynomial pressure dependence of the fragility parameter as well as the zero-mobility temperature captures the evolution of dynamics across broad regions of thermodynamic phase space. The surfaces are constrained by atmospheric-pressure data and by an independently measured isochrone obtained using the \tgp technique, in which temperature is systematically varied under pressure to trace a locus of constant relaxation time near the glass transition, thereby providing an experimental determination of the $T_g(P)$ relation. For glycerol, the framework captures deviations from simple scaling associated with pressure-induced suppression of hydrogen bonding. For the van der Waals liquid cumene, a pressure-modified VTF relation enforces the required high-pressure asymptotic approach to vanishing mobility. Together, these experimental and modeling developments provide a coherent framework for measuring and describing structural relaxation in glass-forming liquids under extreme compression.
}

%% START THE FRONTMATTER
%
\begin{frontmatter}

%% TITLE PAGES
%
%  This command generates the title, copyright, and signature pages.
%
\makefrontmatter 
%% DEDICATION
%
%  You have three choices here:
%    1. Use the ``dedication'' environment. 
%       Put in the text you want, and everything will be formated for 
%       you. You'll get a perfectly respectable dedication page.
%   
%
%    2. Use the ``mydedication'' environment.  If you don't like the
%       formatting of option 1, use this environment and format things
%       however you wish.
%
%    3. If you don't want a dedication, it's not required.
%
%
\begin{dedication} 
  This work is dedicated to the community of the Physics Department at the University of Arkansas, whose culture of curiosity, collaboration, and encouragement has made it an exceptional place for academic growth and discovery. It has been both a privilege and a great source of pride to contribute to this environment, particularly in working alongside and helping develop the next generation of physicists and engineers. The persistence, creativity, and ambition of these students continue to inspire my own work, and I am grateful to be part of a community that values both rigorous inquiry and genuine mentorship.
\end{dedication}


%%%%%%%%%%%%
% Place here the duplication release.
% on second thought - put it into the frontmatter command because it is static - just sign it when printed.
%
%%%%%%%%%%%%




%% EPIGRAPH
%
%  The same choices that applied to the dedication apply here.
%
\begin{epigraph} % The style file will position the text for you.
  \emph{All models are wrong, but some are useful.\\
  --George E. P. Box}\\
\end{epigraph}

% \begin{myepigraph} % You position the text yourself.
%   \vfil
%   \begin{center}
%     {\bf Think! It ain't illegal yet.}
% 
% 	\emph{---George Clinton}
%   \end{center}
% \end{myepigraph}

%\begin{acknowledgements} 
% Thank you to people, things, and places. Without them this thesis would be a meaningless arrangement of ink dots.
%\end{acknowledgements}

\tableofcontents
\listoffigures  % Uncomment if you have any figures
\listoftables   % Uncomment if you have any tables
\newpage


%\begin{vitapage}
%\begin{publications}
%\item{``\textit{Title of published manuscript 1}'', Adequate citation style, \textit{Published}.}
%\item{``\textit{Title of submitted manuscript not yet published}'', Authors, \textit{Submitted}.}
%\item{``\textit{Title of manuscript to be submitted }'', Authors,  (\textit{To be Submitted}).}
%\end{publications}
%\end{vitapage}

\end{frontmatter}

\clearpage
\pagenumbering{arabic}
\setcounter{page}{1}